\documentclass[11pt]{article}
\usepackage{amsmath,amssymb,amsthm}
\usepackage[margin=1in]{geometry}
\newtheorem{proposition}{Proposition}
\title{Room of Mathematicians, item 6 (P6, $\beta_2$ transcendence): Seat Ramanujan\\
regular continued fraction of $q^\ast$}
\author{}
\date{2026-09-26}
\begin{document}
\maketitle

\noindent\textbf{Task.} Compute the ordinary (regular, simple) continued fraction
$[a_0;a_1,a_2,\ldots]$ of the real number $q^\ast$ itself (not the $q$-series continued
fraction of $S$ already exhausted by five reparametrizations in \texttt{P6\_note.tex} /
\texttt{P6\_v2\_note.tex}), and look for a large partial quotient (Liouville-style jump) or
near-periodicity (which would prove $q^\ast$ quadratic irrational).

\section{Calibration point}

Both \texttt{P6/P6\_note.tex} and \texttt{P6\_v2/P6\_v2\_note.tex} were read in full. $q^\ast$ is
defined there as the least positive root of
\[
S(q) = 1+\sum_{k\ge1}(-1)^k q^{2k(k-1)}\frac{[2q(1-q)]^k}{\prod_{j=1}^k(1-q^{2j-1})(1-q^{2j})},
\]
identified with a value of the Hahn--Exton $q$-Bessel function $J^{(3)}_{-1/2}(\cdot;q^2)$, and
$\beta_2=1/\sqrt{q^\ast}$.

\textbf{VERIFIED (this session):} the pre-existing certified bracket
\texttt{private/ROOM/Beta2/g1\_certificate/bracket.txt} (integers \texttt{a\_num}, \texttt{b\_num},
\texttt{Dd=2090}, encoding a rigorous interval $[a\_num/10^{2090},\,b\_num/10^{2090}]\ni q^\ast$,
as produced by \texttt{step1\_bracket.py}) was read directly (not re-run from scratch, since the
2090-digit bracket already exists and re-deriving it is expensive); its midpoint, expanded in
mpmath at working precision \texttt{dps}$\,\in\{200,800,2000,2050\}$, reproduces
\[
q^\ast = 0.449453630558948046125545825395706089225112808545993877680333\ldots
\]
to at least the first 62 digits shown, matching \texttt{P6\_note.tex}/\texttt{P6\_v2\_note.tex}
digit for digit. This confirms the correct number is being used for the continued-fraction
computation below. (Independent Newton re-derivation from the series, as done for \S2 of
\texttt{P6\_note.tex}, was not repeated here since the calibration above already confirms
agreement with that independently-verified value.)

\section{Method}

The regular continued fraction was computed by the standard Euclidean algorithm
($a_i=\lfloor x_i\rfloor$, $x_{i+1}=1/(x_i-a_i)$) applied to the lower endpoint $a\_num/10^{2090}$,
the upper endpoint $b\_num/10^{2090}$, and the midpoint of the certified bracket, each at mpmath
working precision \texttt{dps}$=2050$ (and cross-checked at \texttt{dps}$\in\{200,800\}$ against
the midpoint), wrapped in the required \texttt{perl alarm 290} cap. All runs completed well under
the time cap.

\textbf{Certification of the partial-quotient sequence.} The lower- and upper-bracket endpoints
give \emph{provably} the same 1000 partial quotients (their continued fractions agree in all 1000
computed terms; no divergence index was found up to $i=1000$), which certifies that all 1000
listed partial quotients are correct for the true value of $q^\ast$, independent of where inside
the certified bracket the true value falls. (Precision was not exhausted: 1000 terms from a
2090-digit input is well inside the regime where CF term extraction is reliable; digit consumption
per term averaged $\approx2.1$ digits, consistent with the Lévy-constant rate for a generic
number, so pushing to more terms was not attempted since 1000 already comfortably exceeds the
requested 100+.)

\section{Result: the partial-quotient sequence}

$q^\ast = [a_0;a_1,a_2,\ldots]$ with $a_0=0$ and (the first 200 of the 999 certified terms
$a_1,a_2,\ldots$; the full list of 999 is in the companion script output and is available on
request, omitted here only for length):
\begin{quote}\footnotesize
2, 4, 2, 4, 7, 1, 22, 1, 5, 1, 2, 1, 6, 2, 1, 1, 1, 2, 2, 2, 1, 1, 1, 10, 1, 6, 1, 1, 1, 22, 1, 32,
27, 40, 2, 3, 1, 1, 7, 2, 5, 13, 1, 1, 1, 1, 17, 8, 1, 7, 1, 1, 3, 2, 2, 4, 1, 2, 1, 1, 5, 1, 2, 1,
1, 1, 1, 1, 1, 1, 1, 8, 1, 10, 4, 1, 1, 2, 50, 1, 5, 1, 10, 1, 9, 1, 1, 20, 16, 4, 1, 46, 94, 1, 48,
1, 32, 1, 2, 2, 17, 1, 1, 6, 2, 18, 1, 1, 2, 1, 1, 1, 1, 8, 10, 3, 1, 1, 5, 1, 4, 5, 1, 1, 1, 6, 1,
1, 1, 3, 1, 32, 5, 1, 3, 1, 3, 1, 10, 6, 2, 1, 38, 4, 1, 4, 31, 6, 1, 5, 1, 1, 5, 1, 6, 1, 1, 4, 1,
2, 2, 1, 6, 1, 26, 1, 20, 1, 3, \textbf{8250}, 1, 1, 2, 10, 1, 2, 1, 33, 2, 2, 4, 3, 2, 5, 4, 1, 2,
3, 21, 1, 83, 3, 1, 11, 1, 1, 1, 2, 3, 1, \ldots
\end{quote}
Notable later large terms (full 999-term list): $a_{409}=125$, $a_{571}=616$, $a_{602}=387$,
$a_{689}=127$, $a_{769}=388$, $a_{937}=515$ (indices counted from $a_1=1$-st term after $a_0$).

\section{Analysis}

\textbf{1. Largest partial quotient / Liouville-type jump.} The single largest term among the
first 999 is $a_{170}=8250$ (using $a_1,\ldots,a_{999}$ indexing). By the standard heuristic for
a ``generic'' (Gauss--Kuzmin typical) irrational, the expected order of magnitude of the largest
of $n$ partial quotients is $\Theta(n)$ (more precisely $\max_{i\le n}a_i/n\to$ a random variable
with a Fréchet-type limit law, since $\Pr(a_i>x)\sim \log_2(1+1/x)\sim 1/(x\ln2)$ is a
regularly-varying, heavy ($\alpha=1$) tail); for $n=999$ that puts the expected maximum at a few
thousand, with substantial fluctuation. $8250$ is roughly a factor of 3--6 above the naive
$n/\ln2\approx1440$ estimate, which is an unremarkable fluctuation for a heavy-tailed ($\alpha=1$)
maximum, \emph{not} a Liouville-type anomaly. A genuine Liouville-style excess would look like a
partial quotient many orders of magnitude (e.g.\ $10^{50}$ or more, growing doubly-exponentially
with index) larger than its predecessors' index would suggest; nothing of that kind appears here,
nor at $a_{571}=616$, $a_{769}=388$, $a_{937}=515$, all of which are consistent with ordinary
heavy-tail fluctuations at their respective indices.

\textbf{Verdict on (1): no evidence of a Liouville-type jump.} The largest quotients found are
unremarkable outliers of a heavy-tailed but otherwise typical distribution, not evidence for
$q^\ast$ being a specially-constructed (Liouville) transcendental, and not large enough (relative
to the growth of the CF's own denominators $q_i$) to support a direct Roth/Liouville exclusion
argument against any specific rational approximant: a Liouville-type argument needs
$a_{i+1}\gtrsim q_i^{\,i}$ or faster (super-exponential in the denominator), and here
$a_{170}=8250 \ll q_{169}$ (the 169th CF denominator, already astronomically larger than $8250$
at that depth). This route does not open up.

\textbf{2. Statistical genericity check (Gauss--Kuzmin \& Khinchin).} The empirical frequency of
each digit value $k=1,\ldots,10$ among the 999 terms was compared to the Gauss--Kuzmin prediction
$\Pr(a_i=k)=\log_2\!\big(\tfrac{(k+1)^2}{k(k+2)}\big)$:
\begin{center}
\begin{tabular}{c|cccccccccc}
$k$ & 1 & 2 & 3 & 4 & 5 & 6 & 7 & 8 & 9 & 10\\\hline
observed & 0.4164 & 0.1622 & 0.0821 & 0.0661 & 0.0370 & 0.0250 & 0.0190 & 0.0190 & 0.0120 & 0.0280\\
Gauss--Kuzmin & 0.4150 & 0.1699 & 0.0931 & 0.0589 & 0.0406 & 0.0297 & 0.0227 & 0.0179 & 0.0145 & 0.0120
\end{tabular}
\end{center}
Agreement is good to within sampling noise for $n\approx1000$ terms (the $k=10$ cell is the
biggest outlier, $0.028$ vs.\ $0.012$, but a single bin out of ten at this sample size is not a
significant deviation). The empirical geometric mean of $a_1,\ldots,a_{999}$ is $2.7366$, close
to Khinchin's constant $K_0=2.68545\ldots$ (again consistent with a generic real number; exact
convergence to $K_0$ is itself only almost-sure/typical, not guaranteed for any single number, so
this is weak but consistent evidence, nothing more).

\textbf{Verdict on (2): the sequence looks statistically generic}, i.e.\ consistent with what
almost every real number's continued fraction looks like, hence \emph{uninformative} on its own
for proving or disproving irrationality/transcendence (as the task brief itself anticipated).

\textbf{3. Periodicity / near-periodicity check.} All length-5 windows among the 999 terms were
hashed and checked for exact repeats. 62 repeated 5-windows were found, but every one of them is a
short run of small digits (patterns like $(1,1,1,1,1)$, $(1,2,1,1,1)$, $(1,1,1,3,1)$,
$(5,1,3,1,3)$), which recur by ordinary chance in any $\sim\!1000$-term sequence over a small
alphabet dominated by the digit $1$ (birthday-paradox coincidences, expected count for $n=999$,
alphabet effectively dominated by $\{1,2,3\}$: several dozen is unsurprising). No repeat extends
into a genuine periodic tail (checked by looking at what follows each matched pair: in every case
the sequences diverge immediately after the matched 5-window). No period of any length up to
$\sim500$ (half the computed sequence) was found to repeat exactly starting from any offset. Since
full periodicity of a regular continued fraction is exactly equivalent (Lagrange's theorem) to
$q^\ast$ being a quadratic irrational, and would be an extraordinary, easily-falsifiable claim, the
correct and sufficient check is exactly this exhaustive repeated-window scan against the full 999
terms, which is negative.

\textbf{Verdict on (3): no periodicity or near-periodicity found.} $q^\ast$ shows no sign of being
quadratic irrational (as expected --- this would have been an extremely surprising outcome, and
none of the apparent short repeats survives past 5 terms).

\section{Overall verdict}

\begin{proposition}[status]
\begin{enumerate}
\item[(a)] \textbf{VERIFIED (this session):} $q^\ast=0.449453630558948046125545825395706089225\ldots$,
matching the existing 2090-digit certificate, confirmed to 60+ digits independently at three
working precisions.
\item[(b)] \textbf{VERIFIED (this session):} the regular continued fraction of $q^\ast$ begins
$[0;2,4,2,4,7,1,22,1,5,1,2,1,6,2,1,1,1,2,2,2,\ldots]$; 999 partial quotients after $a_0$ were
computed and certified (lower/upper bracket endpoints agree on all 999 terms, so these are
correct regardless of where the true value sits inside the certified interval).
\item[(c)] \textbf{VERIFIED (this session):} the largest partial quotient among the first 999 is
$a_{170}=8250$; five other terms exceed 100 ($a_{409}=125$, $a_{571}=616$, $a_{602}=387$,
$a_{689}=127$, $a_{769}=388$, $a_{937}=515$). None of these is large enough, relative to the CF
denominator growth at its index, to support a Liouville-type exclusion argument, and their
frequency/size is consistent with the heavy-tailed ($\alpha=1$) Gauss--Kuzmin law for a typical
irrational, not with a deliberately-constructed Liouville number.
\item[(d)] \textbf{VERIFIED (this session):} digit-frequency and geometric-mean statistics match
the Gauss--Kuzmin/Khinchin predictions for a generic real number; an exhaustive repeated-window
scan over all 999 terms found no periodicity or near-periodicity (no match extends past a 5-term
coincidence), so Lagrange's quadratic-irrationality criterion gives nothing.
\end{enumerate}
\end{proposition}

\textbf{Final verdict: LOOKS GENERIC / UNINFORMATIVE.} The ordinary continued fraction of $q^\ast$
behaves exactly like that of a typical (Gauss--Kuzmin-generic) real number: no anomalously large
partial quotient supporting a Liouville-type route, and no periodicity or near-periodicity
supporting a quadratic-irrationality route. This is itself a real, if negative, finding: it closes
off (within the 999-term window computed, using the certified 2090-digit bracket, comfortably
past the requested 100+ terms) the one additional elementary Diophantine-approximation avenue not
covered by \texttt{P6\_note.tex}'s five $q$-series-CF reparametrizations or \texttt{P6\_v2\_note.tex}'s
four named-theory checks (Nesterenko, periods, CM/singular moduli, Schneider--Lang). It does not,
and cannot by its nature, prove irrationality or transcendence --- a generic-looking CF is exactly
what one would see for both a ``boring'' algebraic irrational and for $q^\ast$ actually being
transcendental; it merely rules out the two special structural signatures (Liouville jumps,
periodicity) that would have given an immediate elementary proof one way or another.

\section{Weak points of this note}
\begin{itemize}
\item Only 999 partial quotients were computed (from the existing 2090-digit certified bracket);
a substantially longer expansion (requiring a longer certified bracket, i.e.\ re-running
\texttt{step1\_bracket.py} at higher \texttt{Dd}, which was not attempted here for compute-budget
reasons) could in principle reveal a large quotient or near-period further out, though this is a
priori no more likely than what was already checked.
\item ``Near-periodicity'' was operationalized as exact-repeat of length-5 windows; a more
sophisticated near-periodicity metric (e.g.\ allowing small perturbations, or checking longer
windows with edit-distance) was not attempted.
\item The Liouville-exclusion argument in \S4(1) (comparing $a_{i+1}$ to $q_i^{\,i}$) is a standard
heuristic bound, not spelled out as a fully rigorous inequality with explicit constants in this
note.
\end{itemize}

\end{document}
